\section{Introduction}
\label{sec:introduction}

Requirements Engineering (RE) serves as the foundational cornerstone of the software development lifecycle. Flaws and omissions introduced during early requirements activities propagate exponentially downstream, resulting in severe budgetary overruns, architectural rework, and project abandonment \cite{pohl2016requirements, davis2018maturity}. Despite decades of methodological evolution, eliciting complete requirements and specifying them unambiguously remains one of the most stubborn challenges in software engineering practice.

\subsection{Problem statement and cognitive barriers}
A primary root cause of requirements defects resides in human cognitive and communication limitations. In software engineering, natural language specifications are inherently prone to incompleteness and ambiguity due to cognitive-linguistic transformational effects, such as nominalization, universal quantifiers, and incompletely specified conditions~\cite{pohl2016requirements}. Furthermore, foundational RE frameworks, including the International Requirements Engineering Board (IREB) Certified Professional for Requirements Engineering (CPRE) curriculum \cite{pohl2016requirements}, emphasize that requirements elicitation is an act of mutual knowledge construction between heterogeneous stakeholders. Domain experts, end users, and engineers inhabit disparate conceptual worlds characterized by distinct professional vocabularies, mental models, and unstated assumptions. Stakeholders routinely take implicit background knowledge for granted, omitting critical domain logic from textual Software Requirements Specifications (SRS).

According to IREB Principle 4 (\textit{Adequate Change of Documentation Type}), prototyping mitigates these linguistic vulnerabilities by transforming abstract textual descriptions into concrete behavioral artifacts \cite{pohl2016requirements}. This transformation allows stakeholders to interact directly with the system, surfacing tacit assumptions and uncovering both subconscious baseline expectations (dissatisfiers) and unconscious value-adding needs (delighters) mapped across the Kano model. Furthermore, lengthy textual documents often cultivate an ``illusion of completeness,'' wherein passive approval of textual specifications masks critical omissions, lexical ambiguities, and divergent interpretations \cite{ferrari2016ambiguity, cascon2025ambiguity}.

\subsection{The prototyping paradigm}
Prototyping offers a tangible, experiential mechanism to bridge this cognitive gap. By transforming abstract textual representations into concrete visual and interactive artifacts, prototyping establishes a shared mental model (\textit{shared cognitive artifact}) across technical and non-technical stakeholders \cite{bjarnason2023empirically}. Prototyping facilitates participative evaluation \cite{hickey1998prototyping}, stimulates conversational exploration, and grounds abstract requirements in contextual workflows.

Nevertheless, empirical questions persist regarding whether prototyping universally outperforms conventional elicitation and specification techniques, or whether its efficacy depends upon fidelity, interaction modalities, and specific requirement classifications.

\subsection{Terminology and conceptual foundations}
To establish conceptual clarity throughout this review, core constructs are operationalized following standard RE terminology. \textit{Missing requirements} denote functional rules, exceptional flows, boundary constraints, or system quality attributes that are necessary for fulfilling stakeholder goals but remain entirely omitted from specifications due to tacit knowledge lock-in or inadequate elicitation prompts. \textit{Unclear requirements} refer to requirement statements containing linguistic ambiguities (lexical, syntactic, scope) or non-verifiable subjective adjectives (vagueness), allowing multiple conflicting behavioral interpretations among stakeholders and developers \cite{cascon2025ambiguity}. \textit{Prototyping modality and fidelity} encompass the medium and level of realism embodied in a prototype, spanning Low-Fidelity (Lo-Fi: paper cutouts, sketches, wireframes emphasizing functional flow without visual polish) and High-Fidelity (Hi-Fi: clickable Figma mockups, executable front-ends incorporating interactive behavior and realistic data). Finally, the \textit{Prototyping Assessment Model (PAM)} proposed by Bjarnason et al. \cite{bjarnason2023empirically} serves as the primary multi-dimensional taxonomy, classifying prototyping interventions across five interdependent facets: purpose, scope, media, exploration strategy, and use.

\subsection{Research questions}
To systematically investigate this topic, this review is structured around three core Research Questions:
\begin{itemize}
    \item \textit{RQ6.1:} How are prototypes used to support requirements elicitation and validation?
    \item \textit{RQ6.2:} Does using a prototype help stakeholders identify missing or unclear requirements better than using textual specifications alone?
    \item \textit{RQ6.3:} What types of requirements can be clarified more effectively through prototyping?
\end{itemize}

To enforce analytical rigor, primary studies are evaluated across three structured dimensions: (1) \textit{Rationale (The ``Why'')}, examining what cognitive barriers, tacit knowledge gaps, or textual specification failures necessitate prototyping; (2) \textit{Operational Protocol (The ``How'')}, detailing what empirical elicitation workflows, interaction protocols (e.g., Think-Aloud, Wizard-of-Oz), and validation sessions are executed; and (3) \textit{Taxonomical Classification (The ``What'')}, determining how prototypes are categorized across fidelity levels, media formats, and requirement types under models such as the PAM framework \cite{bjarnason2023empirically}.

The remainder of this paper is structured across seven formal sections: Section~\ref{sec:methodology} delineates the review methodology, database search protocol, and inclusion criteria. Section~\ref{sec:thematic} presents a thematic review addressing RQ6.1 and RQ6.3. Section~\ref{sec:comparative} conducts a comparative analysis addressing RQ6.2. Section~\ref{sec:gaps} analyzes critical research gaps. Section~\ref{sec:directions} outlines proposed research directions. Finally, Section~\ref{sec:conclusion} concludes the paper.
